Modernising ageing safety-critical electronics is often more cost-effective and environmentally sustainable than complete replacement---provided that functional equivalence with the original system can be demonstrated and re-certified according to contemporary, more stringent standards. This case study addresses this challenge through the example of a legacy anti-skid protection module used in European public transportation.

A legacy 8-bit controller board (\gls{mcs48} based) was retrofitted with a low-power, flash-based \gls{fpga} (\texttt{Actel A3P1000}) emulating the original \gls{cpu} and \gls{io} logic, complemented by a 32-bit Arm\textsuperscript{\textregistered}~\gls{mc} (\texttt{STM32F401RET6}) for enhanced data logging. A central aim was to develop a methodological testing framework suitable for broader application in safety-critical system retrofits. The core of this methodology is a two-layer verification strategy:
\begin{enumerate*}[label=(\arabic*)]
	\item \textit{Component \& Interface Tests:} Deterministic \gls{hdl} simulation and hardware tests (boundary-scan, timing) validate logical equivalence and interface timing (e.g., watchdog, frequency sweep).
	\item \textit{System-Level Tests:} Scenario-based, non-deterministic testing on a dedicated fixture~\cite{dolata2023} examines integrated behaviour (fault memory, self-tests, diagnostics).
\end{enumerate*}
Results were mapped to Functional Safety standards including \gls{iec61508}~Ed.~2, \gls{en50129}:2018, \gls{en50155}:2021, \gls{en50716}:2023, and \gls{ieee1012}:2024. Although comprehensive \gls{hil} test rigs were beyond this project's scope, the implemented approach mirrors \gls{hil} coverage criteria, aiming to facilitate future integration.

Lifecycle studies indicate that refurbishing existing rolling-stock electronics reduces operational costs and lowers the carbon footprint compared to manufacturing new components. Our analysis demonstrates that board-level retrofits provide comparable economic and ecological benefits. Successful implementation by major European railway operators further validates the general applicability of this rigorous verification methodology.